% 00-about.tex: what this document is, and how to read its marks.
\section*{About this version}
\addcontentsline{toc}{section}{About this version}

Version 1 of this report (\code{polaris\_project\_report.pdf}, preserved unchanged beside this
file) described Polaris as a database-design project: twelve tables around one credential, an
entity-relationship model, a lifecycle state machine, and the argument that identity data must be
post-quantum from the first day. That description remains the historical record. Since then the
repository has grown from twelve tables to thirty-six, from one
application to a federation protocol, a transparency layer and an institutional exchange fabric,
and from zero machine-checked invariants to 191. Version 2 describes the system as it stands at
tag \code{v9.345} (9 September 2026), and it is written for a reader who has never opened the
repository.

The paper keeps what Version 1 got right: the token-centred geometry of the schema, used here as
the centre of every map, and the lifecycle machine, with one correction, since the trigger Version 1
recommended for enforcing it now exists.
The post-quantum argument is kept, as Appendix~\ref{app:from-v1} summarises. Everything else is
written fresh from the tree.

\subsection*{Four marks}

Every capability carries one of four marks, so that code that runs is never confused with an
intention.

\begin{description}
  \item[\sImpl] A path that runs at \code{v9.345}, exercised by a test or a drill in continuous
    integration, and pinned by an invariant check that fails the build if it stops being true.
  \item[\sExp] A path that runs but is limited in a way the repository itself names: not covered by
    CI, notional in scale, or dependent on a backend that CI simulates.
  \item[\sExt] A property the repository implements and tests internally but that no external
    party has yet audited, attacked, deployed or independently re-implemented: no penetration
    test, cryptographic audit, pilot or third-party implementation has occurred.
  \item[\sFut] A proposal: a roadmap row, a design note, or an idea that exists only in this paper.
    It is not code.
\end{description}

Numbers were measured on the tree at \code{v9.345} by the methods in Appendix~\ref{app:counts},
not copied from earlier documents. Where the repository's own documents disagree with its code, the
code wins and the disagreement is recorded in the author's notes that accompany this version.

\subsection*{The analogy, and its limit}

The paper explains Polaris through one sustained analogy: a walled town. The credential record is
the town hall's register; the constraints are the wall; the endpoints are gates; the transparency
logs are watchtowers; independent witnesses stand outside the wall watching the towers; other
authorities are other towns, joined by roads that run one way; and above the whole town stand laws
that bind even the town itself. The analogy describes arrangement and sight lines, not political
centralisation: Polaris is federated per authority, with no central database, root or verification
service, and every recurrence of the analogy returns to the actual mechanism by its real name.

\subsection*{One thesis}

The system is an attempt to build identity infrastructure in which important claims do not have to
be trusted merely because a central application says they are true. Authority, state, history,
privacy boundaries and the system's own behaviour are progressively exposed to independent
verification and constrained by layers that do not share each other's assumptions. Each layer exists
because the layer inside it could be entirely correct and still fail to prove what a society needs
to know. That progression is the spine of the paper.
